\documentclass[10pt,a4paper]{amsart}
\usepackage[margin=19mm]{geometry}
\usepackage{amsmath,amssymb,mathtools}
\usepackage[scr=rsfs,cal=euler]{mathalfa}
\usepackage{xfrac}
\usepackage[unicode,hidelinks,bookmarks=false]{hyperref}
\newcommand{\ie}{i.e.\ }
\newcommand{\cf}{cf.\ }
\newcommand{\resp}{respectively\ }
\newcommand{\Subsection}{\subsection}
\providecommand{\nobreakdash}{\nobreak\hbox{-}}
\setlength{\emergencystretch}{2em}
\multlinegap0pt
\usepackage{bm}
\usepackage{bbm}
\usepackage{dsfont}
\usepackage{xspace}
\usepackage{cancel}
\usepackage{mathtools}
\usepackage[shortlabels]{enumitem}
\usepackage{amsthm}
\makeatletter
\@ifundefined{thm}{\newtheorem{thm}{Thm}}{}
\@ifundefined{defn}{\newtheorem{defn}[thm]{Definition}}{}
\@ifundefined{lem}{\newtheorem{lem}[thm]{Lemma}}{}
\makeatother
\title[The Antipodes of $q$-Quasi-Symmetric Functions and Non-Commu]{The Antipodes of $q$-Quasi-Symmetric Functions and Non-Commutative Quasi-Symmetric Functions}
\author{Shaul Zemel}
\date{Source: arXiv:2607.07870v1 (CC BY 4.0)}
\begin{document}
\maketitle

\noindent\emph{Source and licence.} The Antipodes of $q$-Quasi-Symmetric Functions and Non-Commutative Quasi-Symmetric Functions. Version arXiv:2607.07870v1 (CC BY 4.0), distributed under the Creative Commons Attribution 4.0 International licence, as recorded in the arXiv licence metadata for this exact version and shown on the abstract page https://arxiv.org/abs/2607.07870v1. Full rights notice: licenses/arxiv-2607.07870-CC-BY-4.0.txt.\par\smallskip

\noindent\textit{Editorial scope.} Counted unit: the Lem whose source label is \texttt{matchsigma} in the article (source0.tex lines 1264--1266, source label \texttt{matchsigma}), delivered together with the article's own complete original proof of it (source0.tex lines 1268--1284, one \texttt{proof} environment of 17 lines). No mathematics is written, repaired or re-derived: statement and proof are the article's own text line for line. Required context retained uncounted: sumFnTNC (lines 1086--1088); shuffle (lines 1218--1228); defsigma (lines 1260--1262). Every definition, display and label of those lines is retained unchanged. Omitted from this package and not redistributed: 1--1085, 1089--1217, 1229--1259, 1263--1263, 1267--1267, 1285--1646. Nothing inside a retained excerpt is omitted or altered: every retained body is delivered exactly as the article's lines are written, so no omission receipt of any kind accompanies this package. Citations: the retained excerpts carry no citation command at all, so no bibliography entry of the article is transcribed and no reference is redistributed. Reference adaptation: none was needed and none was made; every reference inside the delivered unit resolves inside it, so no unresolved reference or citation is produced. Printed slips kept literal: no printed word, symbol, hypothesis or number of the counted statement or of its proof was altered. Layout and class: the article's own class is \texttt{article} with packages including amsmath, amsfonts, amsthm, amssymb, mathabx, enumerate; the wrapper is the lane's pinned \texttt{amsart} editorial wrapper at 10pt on A4, which restates the theorem-like environments the retained text needs and adds the article's own macro definitions that the retained text needs (none), with extra wrapper packages (bm, bbm, dsfont, xspace, cancel, mathtools, enumitem). The delivered numbering is the unit's own. Assets: no retained span contains a figure environment, a graphics-inclusion command, a PDF-page inclusion, a table or a TikZ picture; no image file is copied into this package, the package declares no asset dependency and no image file is redistributed with it. Licence: version arXiv:2607.07870v1 of this article is distributed under the Creative Commons Attribution 4.0 International licence, as recorded in the arXiv licence metadata for this exact version and shown on the abstract page https://arxiv.org/abs/2607.07870v1. A full rights notice is delivered at licenses/arxiv-2607.07870-CC-BY-4.0.txt. Characters: every retained excerpt is pure ASCII, so the delivered engine, XeLaTeX with its default Latin Modern text font, reports no missing character.
\begin{lem}
We have $\mathrm{F}_{n,T}^{\rho}=\sum_{J\in\mathcal{M}_{n,T}}x^{\rho J}$ for every such $n$, $T$, and $\rho$. Equivalently, by setting $\mathcal{M}_{n,T}^{\rho}:=\{\rho J\;|\;J\in\mathcal{M}_{n,T}\}$ we get $\mathrm{F}_{n,T}^{\rho}=\sum_{\eta\in\mathcal{M}_{n,T}^{\rho}}x^{\eta}$, which is the sum over all monomials $x^{\eta}$ of degree $n$ in which $\eta_{\rho(k)}\leq\eta_{\rho(k+1)}$ wherever $1 \leq k<n$, and the equality is strict for any $k \in T$. \label{sumFnTNC}
\end{lem}

\begin{defn}
Let $n$ and $m$ be two non-negative integers.
\begin{enumerate}[$(i)$]
\item A \emph{shuffle of size $(n,m)$} is a sequence $\sigma$ of length $m+n$ which contains $n$ instances of the letter l and $m$ instances of the letter r. We write $\mathcal{S}_{n,m}$ for the set of such shuffles, whose size is clearly $\frac{(n+m)!}{n!m!}$.
\item For $\sigma\in\mathcal{S}_{n,m}$ and $1 \leq t<n+m$, write $t$ as $a+b$ where among the first $t$ entries of $\sigma$ there are $a$ instances of l and $b$ of r. Then the entries in the locations $t$ and $t+1$ of $\sigma$ can be ll, lr, rl, and rr.
\item Assume that $n$ and $m$ are positive, and take $\sigma\in\mathcal{S}_{n,m}$, $T\subseteq\mathbb{N}_{n-1}$ and $S\subseteq\mathbb{N}_{m-1}$. We define $R^{\sigma}(T,S)\subseteq\mathbb{N}_{n+m-1}$ by running over $1 \leq t<n+m$, and consider the cases from part $(ii)$ and the parameters $a$ and $b$, as determined by $\sigma$. In the ll case we put $t$ inside $R^{\sigma}(T,S)$ if and only $a \in T$, and in the rr case we do the same with whether $b \in S$. We always put $t$ in $R^{\sigma}(T,S)$ in the rl case, and never in the lr case.
\item For $\sigma\in\mathcal{S}_{n,m}$ with positive $n$ and $m$, and for compositions $\alpha \vDash n$ and $\beta \vDash m$, we set $T:=\operatorname{comp}_{n}^{-1}\alpha\subseteq\mathbb{N}_{n-1}$ and $S:=\operatorname{comp}_{m}^{-1}\beta\subseteq\mathbb{N}_{m-1}$, and define $\gamma^{\sigma}(\alpha,\beta) \vDash n+m$ to be the image of $R^{\sigma}(T,S)$ from part $(iii)$ under
    $\operatorname{comp}_{n+m}$.
\item Let $\rho \in S_{n}$ and $\varphi \in S_{m}$ be permutations, both written in one-line notation, and take $\sigma\in\mathcal{S}_{n,m}$. Write $\varphi+n$ for the sequence obtained by increasing all the entries of $\varphi$ by $n$. We define $\psi^{\sigma}(\rho,\varphi) \in S_{n+m}$ to be the permutation in $S_{n+m}$ that is obtained by putting $\rho$ in the locations marked with l in $\sigma$ and $\varphi+n$ in those containing r, preserving the orders on both.
\end{enumerate} \label{shuffle}
\end{defn}

\begin{defn}
For positive $n$ and $m$ and permutations $\rho \in S_{n}$ and $\varphi \in S_{m}$, take subsets $T\subseteq\mathbb{N}_{n-1}$ and $S\subseteq\mathbb{N}_{m-1}$, and let $\eta\in\mathcal{M}_{n,T}^{\rho}$ and $\kappa\in\mathcal{M}_{m,S}^{\varphi}$ be elements of the sets from Lemma \ref{sumFnTNC}. Then we define $\sigma(\eta,\kappa)\in\mathcal{S}_{n,m}$ by induction as follows. Assume that the first $a+b$ entries in $\sigma$ were determined, with $a$ of them being l's and the other $b$ are r's (the initial step is the case where $a=b=0$). If $a=n$ or $b=m$ then we complete with the remaining symbol to land inside $\mathcal{S}_{n,m}$. Otherwise, we consider the indices $\eta_{\rho(a+1)}$ and $\kappa_{\varphi(b+1)}$, and we put the symbol l as the next entry of $\sigma$ in case $\eta_{\rho(a+1)}\leq\kappa_{\varphi(b+1)}$, and fill it in with an r when $\eta_{\rho(a+1)}>\kappa_{\varphi(b+1)}$. \label{defsigma}
\end{defn}

\begin{lem}
For $n$, $m$, $T$, and $S$ as in Definition \ref{defsigma}, and take an element $\sigma\in\mathcal{S}_{n,m}$. Then the concatenation $(\eta,\kappa)\mapsto\eta\kappa$ is a bijection between the set of pairs $(\eta,\kappa)\in\mathcal{M}_{n,T}^{\rho}\times\mathcal{M}_{m,S}^{\varphi}$ for which $\sigma(\eta,\kappa)=\sigma$ and the set $\mathcal{M}_{n+m,R}^{\psi}$, with $\psi:=\psi^{\sigma}(\rho,\varphi)$ and $R:=R^{\sigma}(T,S)\subseteq\mathbb{N}_{n+m-1}$ from Definition \ref{shuffle}. \label{matchsigma}
\end{lem}

\begin{proof}
This concatenation takes a pair of sequences $\eta$ and $\kappa$ of lengths $n$ and $m$ to a sequence $\theta:=\eta\kappa$ of length $n+m$, and it is clear that every sequence $\theta$ is a unique such composition, where $\eta$ is the first $n$ elements of $\theta$, and $\kappa$ is the last $m$ entries. We thus need to show that when $\theta=\eta\kappa$ we have $\theta\in\mathcal{M}_{n+m,R}^{\psi}$ if and only if $\eta\in\mathcal{M}_{n,T}^{\rho}$, $\kappa\in\mathcal{M}_{m,S}^{\varphi}$, and $\sigma(\eta,\kappa)=\sigma$.

We first assume that $\theta\in\mathcal{M}_{n+m,R}^{\psi}$, and note that for every $1 \leq j<n$, the index $\rho(j)$ is $\psi(k)$ where $k$ is the $j$'s location of l in $\sigma$. Then $\rho(j+1)$ is $\psi(h)$ for the next location of l in $\sigma$, so that $h>k$ and thus $\psi(h)\geq\psi(k)$. Moreover, the only case where $\theta_{\psi(k)}=\theta_{\psi(h)}$ is where none of the entries between $k$ and $h-1$ are in $R^{\sigma}(T,S)$. This means that there can be no r in between (for rl producing an element of that set), and thus $h=k+1$, and for $R^{\sigma}(T,S)$ not to contain $k$ we need $j \not\in T$.

This proves that $\eta\in\mathcal{M}_{n,T}^{\rho}$, and a similar argument with $\varphi+n$ implies that $\kappa\in\mathcal{M}_{m,S}^{\varphi}$. It also shows that the only situation where $\theta_{\psi(k)}=\theta_{\psi(h)}$ for indices $k$ and $h$ for which $\psi(k)=\varphi(j)+n$ and $\psi(h)=\varphi(j+1)+n$ is where $h=k+1$ and $j \not\in S$. We need to show that $\sigma(\eta,\kappa)$ from Definition \ref{defsigma} is our $\sigma$.

For this, we take an index $a+b+1$, with $a<n$ and $b<m$, and assume that we verified that the first $a+b$ entries of $\sigma(\eta,\kappa)$ coincide with those of $\sigma$, and contained $a$ instances of l and $b$ instances of $r$. Definition \ref{shuffle} shows that $\psi(a+b+1)$ equals $\rho(a+1)$ (resp. $\varphi(b+1)+n$) if that location of $\sigma$ is l (resp. r), and if we define $c>a+b+1$ be the minimal location of an r (resp. l) in $\sigma$ then $\psi(c)$ is $\varphi(b+1)+n$ (resp. $\rho(a+1)$). Since $\theta_{\psi(a+b+1)}\leq\theta_{\psi(c)}$, we deduce in the l case that $\eta_{\rho(a+1)}\leq\kappa_{\varphi(b+1)}$ and indeed the location number $a+b+1$ in $\sigma(\eta,\kappa)$ is also l. In the r case we get $a+b+1 \leq c-1 \in R$ and thus $\kappa_{\varphi(b+1)}=\theta_{\psi(a+b+1)}<\theta_{\psi(c)}=\eta_{\rho(a+1)}$, so that this location is r also in $\sigma(\eta,\kappa)$.

Therefore all the locations in $\sigma(\eta,\kappa)$ coincide with those in $\sigma$ until we cover either all the l's or all the r's. But then the rest is also the same since both of these sequences are in $\mathcal{S}_{n,m}$. We have thus established this direction.

Conversely, assume that $\eta$ and $\kappa$ are in the corresponding sets and produce $\sigma$ via Definition \ref{defsigma}. We need to show that $\theta_{\psi(t)}\leq\theta_{\psi(t+1)}$ for any $1 \leq t<n+m$, with the inequality being strict when $t \in R$. There are the four cases for the locations $t$ and $t+1$ of $\sigma$, namely ll, lr, rl, and rr as in Definition \ref{shuffle}. We can write $t=a+b$, where the first $t$ entries of $\sigma$ consist of $a$ l's and $b$ r's.

But Definition \ref{shuffle} implies that $\theta_{\psi(t)}$ is $\eta_{\rho(a)}$ in the former two cases and $\kappa_{\varphi(b)}$ in the latter two cases, while the value of $\theta_{\psi(t+1)}$ is $\eta_{\rho(a+1)}$ in the first and third cases, and it is $\kappa_{\varphi(b+1)}$ in the second and fourth ones. Moreover, we deduce from Definition \ref{defsigma} that the first and third case correspond to the inequality $\eta_{\rho(a+1)}\leq\kappa_{\varphi(b+1)}$, while the second and fourth one are associated with $\eta_{\rho(a+1)}>\kappa_{\varphi(b+1)}$. Applying these considerations to the $t$th location implies that $\eta_{\rho(a)}\leq\kappa_{\varphi(b+1)}$ in the first two cases, and produces the inequality $\kappa_{\varphi(b)}<\eta_{\rho(a+1)}$ in last two.

Now, our assumption on $\eta$ and $\kappa$ produce $\eta_{\rho(a)}\leq\eta_{\rho(a+1)}$ and $\kappa_{\varphi(b)}\leq\kappa_{\varphi(b+1)}$, which yields $\theta_{\psi(t)}\leq\theta_{\psi(t+1)}$ in all four cases. As $t \not\in R$ in the lr case, and we have a strict inequality in rl case, these two cases are covered. As in the first (resp. fourth) case we have $t \in R$ if and only if $a \in T$ (resp. $b \in S$), which then implies $\eta_{\rho(a)}<\eta_{\rho(a+1)}$ (resp. $\kappa_{\varphi(b)}<\kappa_{\varphi(b+1)}$), we deduce the required strict inequality also in these cases. Hence $\theta$ satisfies all the conditions for being in $\mathcal{M}_{n+m,R}^{\psi}$, as desired. This completes the proof of the lemma.
\end{proof}

\end{document}
